% !TEX root = ./fduthesis.tex

% Load the template fonts before xeCJK checks the default CJK family.
\DeclareHookRule{begindocument/before}{fduthesis}{before}{xeCJK}

% setup.tex
\fdusetup{
  style = {
    font = times,
    cjk-font = windows,
    font-size = -4,
    fullwidth-stop = false,
    hyperlink = none,
    auto-make-cover = false, % 暂隐封面、指导小组成员页与封三声明页；恢复时改回 true 即可
    bib-backend = biblatex,
    bib-resource = {citelist.bib},
    bib-style = gb7714-2025
  },
  info = {
    degree = academic,
    title = {Hirota--Maxwell--Bloch 方程的非零边界谱分析与孤子解},
    title* = {Spectral Analysis and Soliton Solutions of the Hirota--Maxwell--Bloch System with Nonzero Boundary Conditions},
    author = {某\quad 某},
    author* = {Mou Mou},
    supervisor = {某\quad 某\quad 教~授},
    department = {数学科学学院},
    major = {应用数学},
    student-id = {202XXXXXXXX},
    instructors = {{某\quad 某\quad 教~授},{某\quad 某\quad 教~授},{某某某\quad 教~授}},
    keywords = {Hirota--Maxwell--Bloch 系统, 非零边界条件, 逆散射变换, Riemann--Hilbert 问题, 介质重构, 孤子解},
    keywords* = {Hirota--Maxwell--Bloch system, Riemann--Hilbert problem, nonzero boundary conditions, inverse scattering transform, medium reconstruction, soliton solutions},
    clc = {O175.29}
  }
}

% 字体，英文字体更新为Times New Roman
\setmainfont{Times New Roman}

% 序号
\usepackage{enumerate}

% 图片与浮动控制
\usepackage{graphicx,placeins}

% 子图宏包
\usepackage{subcaption}

% 表格排列与三线表
\usepackage{array}
\usepackage{booktabs}
\usepackage{longtable}

% 代码环境宏包（更新区需要）
\usepackage{listings}
\lstdefinestyle{lstStyleCode}{%
  aboveskip         = \medskipamount,
  belowskip         = \medskipamount,
  basicstyle        = \ttfamily\zihao{6},
  commentstyle      = \slshape\color{black!60},
  stringstyle       = \color{green!40!black!100},
  keywordstyle      = \bfseries\color{blue!50!black},
  extendedchars     = false,
  upquote           = true,
  tabsize           = 2,
  showstringspaces  = false,
  xleftmargin       = 1em,
  xrightmargin      = 1em,
  breaklines        = false,
  framexleftmargin  = 1em,
  framexrightmargin = 1em,
  columns           = flexible,
  keepspaces        = true,
  texcl             = true,
  mathescape        = true
}

% 超链接
\usepackage{hyperref}

%更新区------------------------------------------------------------------------
\lstnewenvironment{codeblock}[1][]{%
  \lstset{style=lstStyleCode,#1}}{}
%封面字体更新
\AtBeginDocument{
 \addtocontents{toc}{\protect\hypersetup{linkcolor=black}}
  \addtocontents{toc}{\protect\color{black}}

  \addtocontents{lof}{\protect\hypersetup{linkcolor=black}}
  \addtocontents{lof}{\protect\color{black}}

  \addtocontents{lot}{\protect\hypersetup{linkcolor=black}}
  \addtocontents{lot}{\protect\color{black}}

  \addtocontents{loa}{\protect\hypersetup{linkcolor=black}}
  \addtocontents{loa}{\protect\color{black}}

  \addtocontents{lol}{\protect\hypersetup{linkcolor=black}}
  \addtocontents{lol}{\protect\color{black}}
}
%摘要更新
\ExplSyntaxOn
\cs_set_protected:Npn \__fdu_abstract_begin:
  {
    \cleardoublepage
    \thispagestyle { plain }
    % —— 补回被丢掉的目录项 + PDF 书签 ——
    \phantomsection
    \addcontentsline { toc } { chapter } { \c__fdu_name_abstract_tl }
    % —— 补回页眉 ——
    \__fdu_chapter_header:n { \c__fdu_name_abstract_tl }
    % —— 你的自定义排版 ——
    \zihao { -4 }
    \begin { center }
      \l__fdu_info_title_tl
    \end { center }
    \noindent \c__fdu_name_abstract_tl \c__fdu_fwid_colon_tl
  }
\cs_set_protected:Npn \__fdu_abstract_en_begin:
  {
    \cleardoublepage
    \thispagestyle { plain }
    \phantomsection
    \addcontentsline { toc } { chapter } { \c__fdu_name_abstract_en_tl }
    \__fdu_chapter_header:n { \c__fdu_name_abstract_en_tl }
    \zihao { -4 }
    \begin { center }
      \l__fdu_info_title_en_tl
    \end { center }
    \noindent \c__fdu_name_abstract_en_tl :
  }
\ExplSyntaxOff

\ExplSyntaxOn
\cs_set_protected:Npn \__fdu_abstract_end:
  {
    \__fdu_keywords:nNn
      { \sffamily \c__fdu_name_keywords_tl \c__fdu_fwid_colon_tl }
      \l__fdu_info_keywords_clist { \c__fdu_fwid_semicolon_tl }
  }
\ExplSyntaxOff

\ExplSyntaxOn
\cs_set_protected:Npn \__fdu_abstract_en_end:
  {
    \__fdu_keywords:nNn
      { \bfseries \c__fdu_name_keywords_en_tl \__fdu_quad: }
      \l__fdu_info_keywords_en_clist { ; ~ }
  }
\ExplSyntaxOff
%封面更新
\ExplSyntaxOn
\cs_set_protected:Npn \__fdu_cover_id:
  {
    \noindent
    \begin{minipage}[t]{0.5\textwidth}
      \raggedright
      中图分类号：\l__fdu_info_clc_tl
    \end{minipage}%
    \begin{minipage}[t]{0.5\textwidth}
      \raggedleft
      学校代码：10246 \\
      % 同等学力人员/英文项目   % ← 需要时去掉本行行首的 % 即可显示
    \end{minipage}
  }
\ExplSyntaxOff

\ExplSyntaxOn
\cs_set_protected:Npn \__fdu_cover_info:
  {
    \begin{center}
    \zihao{-3}        % ← 封面字号
    \begin{tabular}{ @{} r @{：} c @{} }
      \makebox[7em][s]{学生姓名} & \l__fdu_info_author_tl       \\[1ex]
      \makebox[7em][s]{学号}     & \l__fdu_info_student_id_tl   \\[1ex]
      \makebox[7em][s]{导师姓名、职称} & \l__fdu_info_supervisor_tl \\[1ex]
      \makebox[7em][s]{学科专业} & \l__fdu_info_major_tl        \\[1ex]
      \makebox[7em][s]{培养单位} & \l__fdu_info_department_tl   \\[1ex]
      \makebox[7em][s]{论文提交日期} & \quad \the\year 年 \the\month 月 \the\day 日 \\
    \end{tabular}
    \end{center}
  }
\ExplSyntaxOff

% 论文正文定理环境与自定义数学记号
\hypersetup{allcolors=black}
\newtheorem{proposition}{命题}
\newtheorem{remark}{注}
\newcommand{\ii}{\mathrm{i}}
\newcommand{\dd}{\mathrm{d}}
\newcommand{\sg}{\sigma_3}
\newcommand{\Id}{I_2}
\newcommand{\conj}[1]{{#1}^{*}}
\DeclareRobustCommand{\Res}{\operatorname*{Res}\limits}
\DeclareMathOperator{\diag}{diag}
\DeclareMathOperator{\tr}{tr}
\DeclareMathOperator{\sech}{sech}
\allowdisplaybreaks[2]
\setlength{\emergencystretch}{2em}
\makeatletter
\setlength{\@fptop}{0pt}
\setlength{\@fpsep}{14pt}
\setlength{\@fpbot}{0pt plus 1fil}
\makeatother
\renewcommand{\topfraction}{0.9}
\renewcommand{\bottomfraction}{0.85}
\renewcommand{\textfraction}{0.1}
\renewcommand{\floatpagefraction}{0.75}
\ctexset{
  figurename = {图},
  listfigurename = {插图清单}
}
